\documentclass[a4paper]{article}
\usepackage{amsmath}
\usepackage{amssymb}
\usepackage{geometry}
\usepackage{hyperref}

\geometry{margin=25mm}

\title{Galactic Rotation Curves without Dark Matter:\\
Compactness of the Dual Spacetime}

\author{https://github.com/hypernumbernet}
\date{\today}

\begin{document}

\maketitle

\begin{abstract}
The same dual-rotor geometry that layers nuclei and electron shells unfolds, at cosmological scales, as compactness. The dual rotor is a three-sphere. We read the gravitational channel, driven by $(\phi,\theta)$ and tied to inertial mass, as organizing that sphere into galaxy-scale patches of compactification radius $R$; the electromagnetic channel, driven by $(\alpha,\beta)$ and tied to charge, unfolds over substantially larger, near-cosmic scales. A spherical source on the three-sphere carries a radially harmonic cotangent potential, normalized by Gauss's law. Its acceleration is a sum over images of the source repeated with period $\pi R$. If those images are laid along the usual radius, the circular speed has no finite value at the antipode $r=\pi R$, which for a Milky-Way mass lies within the measured rotation curve. That placement makes the usual radius periodic, while compactness belongs to the dual spacetime alone. Laid instead along the dual direction, the same images give the acceleration $(GM/(rR))\coth(r/R)$. If the compactification radius scales as $R=\sqrt{GM/a_0}$, this is a single relation between the observed and the Newtonian acceleration, and the rotation curve is the baryonic Newtonian curve multiplied by $x\coth x$, with $x^2=a_0r/v_{\text{bar}}^2$. It rises from the centre, turns over, and flattens without a further rise, at a speed that obeys $v^4=GMa_0$ exactly. If galaxies are such patches, the missing acceleration is this geometry, read from baryons alone. That identification is the proposed reading. The acceleration $a_0$ below is an external input, not a consequence of the dual-rotor algebra.
\end{abstract}

\section{Introduction}
\label{sec:introduction}

Double spacetime theory refuses a shared continuum. Each massive particle carries a usual Minkowski frame and a dual frame, and gravity is the torsional mismatch of that pair~\cite{dst-main}. The dual rotor is a three-sphere. At galactic scales the gravitational channel, driven by $(\phi,\theta)$ and tied to inertial mass, is read as organizing that sphere into patches of compactification radius $R$. The electromagnetic channel, driven by $(\alpha,\beta)$ and tied to charge, unfolds on a larger radius. The acceleration $a_0$ that sets $R=\sqrt{GM/a_0}$, and the Hubble interval later read as the electromagnetic radius, are external inputs. Neither is a consequence of the dual-rotor algebra.

Section~\ref{sec:compact} records the three-sphere and its equidistant chart. Section~\ref{sec:cotangent} writes the cotangent potential. Section~\ref{sec:radial} lays the images along the usual radius, and that placement is not the galactic law. Section~\ref{sec:dual} lays the same images along the dual direction. Section~\ref{subsec:btfr} scales the patch with the mass and obtains the acceleration relation. Section~\ref{sec:rotation} builds the rotation curve from the baryonic Newtonian curve. Section~\ref{sec:sky} isolates what the sky can decide.

\section{The compact dual spacetime}
\label{sec:compact}

The dual rotor $R_{\text{dual}} \in \mathrm{SU}(2) \cong S^3$ parameterizes the dual degrees of freedom intrinsic to every massive particle. Its three rotation angles are coordinates on a three-sphere of radius $R$, while the usual spacetime remains flat. Those angles are a coordinate reading. Each particle carries an element of this three-sphere, and the position of particle $b$ relative to particle $a$ is the rotor $\rho(a,b)=a^{-1}b$. It is the identity precisely when the two elements coincide, and the cosine of the angle between them is the real part of $\rho$. A rotation of the three-sphere is a pair of unit quaternions, applied to every element at once. The left factor leaves $\rho$ unchanged. The right factor replaces $\rho$ by $h^{-1}\rho h$, and the angle is unchanged. Moving the pole of the equidistant map to $g$ replaces the coordinate of an element $p$ by the logarithm of $g^{-1}p$. The coordinates change and $\rho$ does not. With the pole placed at $a$, the coordinate of $b$ is the logarithm of $\rho(a,b)$. This rotor is the relation of two particles. It is not the mismatch of the usual rotor and the dual rotor of one particle. The logarithm
\[
\log : S^3 \to \mathfrak{su}(2) \cong \mathbb{R}^3
\]
sends each rotor to a point of a Euclidean ball of radius $\pi R$. It is the three-dimensional azimuthal equidistant projection: faithful at the origin, and singular at the antipode, which is spread along the boundary of the ball. The volume distortion is the Jacobian
\[
\mathcal{J}(\chi) = \frac{\sin^2 \chi}{\chi^2}, \qquad \chi = r/R.
\]
As $\chi \to 0^+$, $\mathcal{J}$ tends to $1$. At the antipode, $\mathcal{J}$ vanishes. A Newtonian calculation assumes $\mathcal{J}=1$ at every radius. On a galactic patch that assumption holds only near the centre. Every direction at geodesic distance $\pi$ lifts to one and the same element. Relative to the pole its rotor is $-1$, and the angle is $\pi$. The angle $\pi$ occurs at this antipode alone. The chart writes that element at the origin, where it also writes the pole. The element $-1$ lies in the centre of the rotation group, and every rotation of the sphere fixes this relative rotor. A generic relative rotor is conjugated by the right factor. That conjugation leaves $-1$ fixed. Where the chart lays the pole and the antipode on one origin, the rotor is what keeps the position. This element $-1$ is the value of the exponential of $J_\varepsilon$ at the parameter $\pi/\varepsilon$. Off the origin the Jacobian $\mathcal{J}$ is the square of $\sin\chi/\chi$, so it vanishes at the antipode. Along $\chi=\varepsilon t$, as $\varepsilon\to 0$, that factor tends to $1$ at every finite $t$, and the exponential tends to the null rotation of the parent theory~\cite{dst-main}. The antipode recedes, and the direction the flat chart does not keep is the null line.

\section{The cotangent potential}
\label{sec:cotangent}

In geodesic polar coordinates on the three-sphere,
\[
ds^2 = dr^2 + R^2 \sin^2(r/R)\,(d\vartheta^2 + \sin^2\!\vartheta\, d\varphi^2),
\]
the area of a geodesic sphere is $4\pi R^2\sin^2(r/R)$. The cotangent
\[
\Phi(r) = -\frac{GM}{R}\cot\!\left(\frac{r}{R}\right)
\]
is the radial potential whose flux through every such sphere equals $-4\pi GM$. Away from the origin the product of the area factor with $\partial_r\Phi$ is the constant $GM$, and the radial Laplacian therefore annihilates $\Phi$:
\[
\partial_r^2\Phi + \frac{2}{R}\cot\!\left(\frac{r}{R}\right)\,\partial_r\Phi = 0.
\]
Differentiating gives the acceleration
\[
g(r) = -\frac{d\Phi}{dr} = -\frac{GM}{R^2\sin^2(r/R)}.
\]
On the open interval $0<r<\pi R$ this acceleration is strictly negative. The force points toward the mass on both sides of the equator. The potential itself is negative if and only if $r<\pi R/2$, and it vanishes only at the equator. The equator is the edge of the negative well. It is not a reversal of the force. Past it, $\Phi$ is positive and still slopes down toward the source; as the geodesic spheres shrink toward the antipode, the same flux is squeezed into a smaller area and $|g|$ grows, still inward.

For $0<r\le R$ the germ at the origin is Newtonian. The potential differs from $-GM/r + GMr/(3R^2)$ by a remainder of order $r^3/R^4$. The correction $GMr/(3R^2)$ is the curvature of the three-sphere. It disappears in the decompactification $R\to\infty$, and with it the whole departure from the Newtonian potential.

\section{Images along the radius}
\label{sec:radial}

The acceleration is a lattice sum. The partial-fraction expansion of the cotangent gives
\[
\frac{GM}{R^2\sin^2(r/R)}=\sum_{n\in\mathbb{Z}}\frac{GM}{(r-n\pi R)^2}.
\]
The term $n=0$ is Newton's; the others are the source seen again at every multiple of $\pi R$ along the radius. A circular orbit in the flat chart has
\begin{equation}
\label{eq:s3-circular}
v^2(r) = r\,|g(r)| = \frac{GM}{R}\,f(r/R), \qquad f(x)=\frac{x}{\sin^2 x}.
\end{equation}
On $(0,\pi/2)$ the shape has a single shallow minimum, at the root $x_0\in(1.14,1.20)$ of $\tan x=2x$, with $1.35<f(x_0)<1.41$. Past it $f$ rises, and as $x\to\pi^-$ it tends to $+\infty$: the image $n=1$ reaches the orbit at the antipode. For a Milky-Way mass, under the scaling of Section~\ref{subsec:btfr}, the antipode lies between $25$ and $29\,\mathrm{kpc}$, and measured rotation curves stay flat across such radii. The radial lattice makes the usual radius itself periodic, whereas only the dual spacetime is compact. It is the chart's reading of the images, faithful near the centre, and it is not the galactic law.

\section{Images along the dual direction}
\label{sec:dual}

The duality $X\mapsto Xi$ relates the elliptic geometry of the dual sector to the hyperbolic geometry of the usual sector by analytic continuation of the parent theory~\cite{dst-main}; on the cotangent it acts as $i\cot(ix)=\coth x$. Carried through that continuation, the period $\pi R$ leaves the usual radius and runs along the dual direction. Image $n$ sits at squared distance $r^2+(n\pi R)^2$, one leg in usual space and one around the compact dual circle, and the usual radius stays non-compact. The same cotangent expansion, at imaginary argument, sums this lattice in closed form:
\begin{equation}
\label{eq:dual-lattice}
g(r)=\sum_{n\in\mathbb{Z}}\frac{GM}{r^2+(n\pi R)^2}=\frac{GM}{rR}\coth\frac{r}{R}.
\end{equation}
The ratio to Newton is the dual enhancement
\[
\eta(x)=x\coth x, \qquad x=r/R,
\]
and for every $x>0$
\[
\max(1,x)<\eta(x)<\min\!\left(1+x,\;1+\frac{x^2}{3}\right).
\]
The enhancement increases strictly and tends to $1$ at the origin, while $\coth$ decreases strictly to $1$. Near the centre the two lattices agree: both begin as $GM/r^2+GM/(3R^2)$, the curvature correction of the cotangent germ. The dual sum stays strictly below that germ at every radius, and strictly above both $GM/r^2$ and $GM/(rR)$. At large radius the second dominates. The acceleration falls as $1/r$ rather than $1/r^2$, and no radius is singular.

At a fixed compactification radius the lattice is linear in the source mass: the acceleration of a sum of masses is the sum of the accelerations, and scaling the mass scales the acceleration by the same factor. The fractional excess over the inverse square depends only on $r/R$. It is positive, and strictly less than $(r/R)^2/3$. On $(0,\pi)$ the three-sphere factor $(x/\sin x)^2$, the reciprocal of the volume Jacobian, lies strictly above the shared germ $1+x^2/3$, while the dual enhancement lies strictly below it. A solar mass, on the same external inputs, has compactification radius between $7020$ and $7040$ astronomical units, so at one astronomical unit the dual fractional excess lies below $10^{-8}$. At the Milky-Way compactification radius the exterior rapidity exceeds its Newtonian value by a fraction below $10^{-6}$, beside a dual excess at $r=R$ between $0.31$ and $0.32$.

\section{Baryonic Tully--Fisher and the acceleration relation}
\label{subsec:btfr}

Let the compactification radius scale with the mass as $R=\sqrt{GM/a_0}$, where $a_0 \approx 1.2 \times 10^{-10}\;\mathrm{m/s^2}$ is the observed baryonic Tully--Fisher acceleration, taken as an input rather than derived from the dual-rotor algebra. A heavier galaxy is a larger patch. A derivation of $a_0$ is not claimed. About a point mass the circular speed is $v^2=r\,g=(GM/R)\coth(r/R)$. It decreases strictly with radius toward
\[
v_\infty^2=\frac{GM}{R}, \qquad v_\infty^4=GM\,a_0,
\]
exactly: slope four, with no shape factor. The approach is from above and fast. Past $r=3R$, $v^2$ lies within $1/200$ of $v_\infty^2$, so $v$ is within a quarter of a percent of $v_\infty$. There is no second rise and no antipode.

The same scaling turns the lattice into a radial acceleration relation. With $g_{\text{N}}=GM/r^2$ one has $x^2g_{\text{N}}=a_0$ and $GM/(rR)=\sqrt{a_0g_{\text{N}}}$, so the acceleration is a function of $g_{\text{N}}$ alone,
\begin{equation}
\label{eq:dual-rar}
g=\sqrt{a_0\,g_{\text{N}}}\,\coth\sqrt{\frac{a_0}{g_{\text{N}}}},
\end{equation}
bounded by
\[
\max\!\left(g_{\text{N}},\sqrt{a_0g_{\text{N}}}\right)<g<g_{\text{N}}+\min\!\left(\sqrt{a_0g_{\text{N}}},\,\frac{a_0}{3}\right).
\]
At high acceleration the excess over Newton stays below $a_0/3$. Tying the patch radius to the enclosed mass makes the relation strictly subhomogeneous and strictly subadditive: for $k>1$, scaling $g_{\text{N}}$ by $k$ multiplies $g$ by a factor strictly smaller than $k$, and the acceleration of two Newtonian contributions taken together is strictly smaller than the sum of the accelerations. While $g_{\text{N}}$ is large compared with $a_0$, the inverse-square reading carries the acceleration, and the fractional excess becomes of order one as $g_{\text{N}}$ passes through $a_0$. At low acceleration $g$ approaches $\sqrt{a_0g_{\text{N}}}$, the deep regime of Milgrom's phenomenology. In between, the relation is fixed, with no interpolating function to choose. At $g_{\text{N}}=a_0$, where $x=1$, the enhancement is $\coth 1$, between $1.31$ and $1.32$.

For the Milky-Way stand-in $M \approx 6 \times 10^{10}\,M_\odot$ the same inputs give $8\,\mathrm{kpc}<R<9\,\mathrm{kpc}$ and
\[
175\,\mathrm{km/s}<v_\infty<176\,\mathrm{km/s}.
\]
To identify $a_0$ with a characteristic acceleration of the patch is still a proposal.

\section{The rotation curve}
\label{sec:rotation}

At radius $r$ the flux is that of the enclosed mass $M(r)$, and the patch that carries it has $R(r)^2=GM(r)/a_0$. Outside the baryons this is the single patch of the total baryonic mass $M_{\text{b}}$; inside, it shrinks with the enclosed mass. With $g_{\text{N}}=v_{\text{bar}}^2/r$, the orbit sits at $x=r/R(r)=\sqrt{a_0/g_{\text{N}}}$, and the rotation curve is
\begin{equation}
\label{eq:dual-rotation}
v^2(r)=v_{\text{bar}}^2(r)\,x\coth x, \qquad x^2=\frac{a_0\,r}{v_{\text{bar}}^2(r)},
\end{equation}
which is $r$ times the relation \eqref{eq:dual-rar} at $g_{\text{N}}=v_{\text{bar}}^2/r$. Here $v_{\text{bar}}^2=v_\star^2+v_{\text{gas}}^2+v_{\text{bulge}}^2$ is computed from photometry and 21-cm maps exactly as in a Newtonian mass model. The stellar mass-to-light ratio is the only adjustable quantity; $a_0$ is universal, and there is no halo. For a spherical mass the recipe is exact. For a disk it applies the relation radius by radius to the Newtonian curve of the disk. The two cases meet in the growth of a galaxy. It begins as an aggregate of stars moving at random. Those angular momenta cancel, and the mass is a ball. The pull follows a great circle of the three-sphere, and the motion is the straight chord of the flat chart, the azimuthal equidistant projection of that sphere. The chart is treated as flat only because a galaxy is small beside the three-sphere of cosmic size; strictly, the motion lies on that huge sphere as well. Place the body at chart distance $\chi$ and the source at chart distance $\psi$, with angle $\alpha$ between those radii. The initial tangent of the great circle, read in the chart, has moment about the source
\begin{equation}
\label{eq:sphere-chord}
\left(\frac{\chi^{2}\sin\psi}{\sin\chi}-\frac{\chi\psi\sin\psi}{\sin\chi}\cos\alpha-\psi\sin\chi\cos\psi+\psi\sin\psi\cos\alpha\cos\chi\right)\sin\alpha,
\end{equation}
normal to the plane of the two radii. The moment vanishes when the body and the source lie on one ray from the chart origin. It vanishes when the source is the origin itself: the great circle through the origin is a straight ray, and the pull stays central in the flat chart. The moment is a reading in that map. A rotation that carries the body, the source, and the pole together leaves the moment unchanged. For one pair the moment is positive when the pole is the identity, and it vanishes when the pole is moved onto the source and the axes are turned so that the body lies along the first coordinate axis. The angle of $\rho$ agrees in the two readings. When the two chart distances are equal and the angle lies in $(0,\pi)$, the moment equals $(1-\cos\alpha)\,\chi\,(\chi-\sin\chi\cos\chi)\sin\alpha$ and is positive. A body moving along the flat chord toward such a source therefore acquires angular momentum about the source. Outside that ball the speed is the point-mass law of the whole baryonic mass, whatever the arrangement inside. Gravity closes on the galactic patch of radius \(R\). The electromagnetic channel reads the same lattice on a larger compactification \(\Lambda\). \(\Lambda\) is an input, as \(a_0\) is. Inside the patch the gravitational acceleration exceeds its inverse-square value by the fraction \(x\coth x-1\), which lies between \(0.31\) and \(0.32\) at \(r=R\). Throughout \(r\le R\) the electromagnetic fraction stays below \((R/\Lambda)^2/3\), and at \(r=R\) the ratio of the gravitational fraction to the electromagnetic one exceeds \(\frac{93}{100}(\Lambda/R)^2\). The separation is a centripetal residual of the gravitational channel, radial about the galactic centre.

The three-sphere admits parallel lines. The pure unit imaginaries determine, at every point, three tangent directions of norm one, orthogonal to one another. Along the left translation \(t\mapsto q(\cos t+I\sin t)\) the displacement rotor is \(\cos s+I\sin s\), independent of the base point and of the instant. Along the right translation \(t\mapsto(\cos t+I\sin t)q\) the same displacement stands on the left. There are therefore two families of parallels, and each passes through every point. Conjugation sends a left line to the right line of the conjugate point and reverses the parameter, so the two families are mirrors. The frames agree at the identity. At the point \(\cos\varphi+I\sin\varphi\) the right-hand \(J\) direction coincides with the left-hand direction \(\cos 2\varphi\, J-\sin 2\varphi\, K\), and for \(0\le\varphi\le\pi/2\) the angle between the frames is \(2\varphi\). Reversing the geodesic angle preserves the cosine and changes only the sign of the \(K\) component. At angle \(\pi/4\) the two \(J\) directions do not agree. This double angle is the twist that the twistor correspondence reads in the Hopf fibration. On a sphere of radius \(\rho\) the mismatch at physical distance \(s\) is \(2s/\rho\). The gravitational channel reads it at radius \(R\), and the electromagnetic channel at radius \(\Lambda\). At one distance the ratio of the mismatches is \(\Lambda/R\), and their difference is \(2s(\Lambda-R)/(R\Lambda)\). At the edge \(s=R\) of the patch the gravitational mismatch is \(2\). Reading the electromagnetic radius as the Hubble length \(c/H_0\), with \(H_0\) an external interval from \(67\) to \(74\,\mathrm{km\,s^{-1}\,Mpc^{-1}}\) containing both the Planck and the distance-ladder central values, and with the Milky Way compactification radius from \(8\) to \(9\,\mathrm{kpc}\), the electromagnetic mismatch at that edge lies between \(3.57\times 10^{-6}\) and \(4.45\times 10^{-6}\), and the gravitational mismatch exceeds it by a factor between \(450138\) and \(559315\). The ratio of the patch radius to the Hubble length therefore lies between \(1.787\times 10^{-6}\) and \(2.222\times 10^{-6}\). Across the patch the equidistant Jacobian on the three-sphere of radius \(\Lambda\) differs from \(1\) by less than \(1.646\times 10^{-12}\), and the electromagnetic fraction stays below the same figure; at the edge the gravitational fraction exceeds the electromagnetic one by more than \(1.8836\times 10^{11}\). As with \(a_0\), the value of \(H_0\) is an input, and the identification is not a consequence of the dual-rotor algebra. Beside the fractional forces, which separate by \((\Lambda/R)^2\), the two parallelizations themselves separate by \(\Lambda/R\). A rotation \(p\mapsto apb\) carries each family onto itself. The right factor conjugates the axis of a left line, and the left factor conjugates the axis of a right line. The mirror exchanges the families, and with them the orientation and the parameter. Through the point \(J\) a quarter turn along the left family ends at \(-K\), and a quarter turn along the right family ends at \(K\). For \(0<\varphi<\pi/2\) the \(K\) component of the frame is \(-\sin 2\varphi\) on the first screw and \(+\sin 2\varphi\) on the reversed geodesic. Both are nonzero, the two are opposite, and a rotation leaves each family on itself. The three-sphere carries one family. Which family it carries is the orientation of this universe, and the mirror universe carries the other.

A source at the chart origin adds none of the moment \eqref{eq:sphere-chord}. A source off the origin adds that moment, in the plane spanned by the origin and the two positions. On the residual, the specific angular momentum of a circular orbit increases strictly with radius and tends to zero at the origin, so each value belongs to exactly one radius. Near the centre a small angular momentum still belongs to a small orbit. Past three compactification radii the speed has levelled, and the specific angular momentum grows with the radius. The moment \eqref{eq:sphere-chord} is odd in the oriented angle between the two radii, so a source and its reflection across the ray sum to zero. Where the chart distances agree and lie in \((0,\pi)\), the two moments are nonzero and opposite. Every patch of this sphere is read against the same screw. The sense of a disk is a sense relative to that screw, one reference at every galaxy, and the mirror universe reverses the reference together with the moment. As the torque carries angular momentum outward, the aggregate expands and keeps a rotation near \(v_\infty\). That rotation is the vortex, and the ball spreads into a disk.

At every radius
\[
\max\!\left(v_{\text{bar}}^2,\sqrt{a_0\,r\,v_{\text{bar}}^2}\right)<v^2<v_{\text{bar}}^2+\sqrt{a_0\,r\,v_{\text{bar}}^2}.
\]
The curve therefore vanishes wherever $r\,v_{\text{bar}}^2$ does, as at a regular centre. Every baryonic curve with a Keplerian tail, $r\,v_{\text{bar}}^2\to GM_{\text{b}}$, flattens at $v_\infty^2=\sqrt{GM_{\text{b}}a_0}$, and outside the baryons the point-mass behaviour above holds with $R=\sqrt{GM_{\text{b}}/a_0}$. The curve rises, turns over, and stays flat.

For a Milky-Way exponential disk of $6\times10^{10}\,M_\odot$ and scale length $2.6\,\mathrm{kpc}$, Eq.~\eqref{eq:dual-rotation} gives about $202\,\mathrm{km/s}$ at $4\,\mathrm{kpc}$, $211$ at $8$, $190$ at $15$, $179$ at $25$, and $176$ at $60\,\mathrm{kpc}$, while the Newtonian disk alone falls from $188$ to $105\,\mathrm{km/s}$ between $8$ and $25\,\mathrm{kpc}$. A low-surface-brightness disk of $3\times10^{9}\,M_\odot$ and scale length $3\,\mathrm{kpc}$, with $g_{\text{N}}$ near $a_0/10$ across the inner disk, rises to about $50\,\mathrm{km/s}$ at $2\,\mathrm{kpc}$, $78$ at $6$ and $87$ at $15\,\mathrm{kpc}$, and settles between $83.1$ and $83.2\,\mathrm{km/s}$. A dwarf of $4\times10^{8}\,M_\odot$ whose baryons lie inside a ball of radius $0.8\,\mathrm{kpc}$ has compactification radius between $0.6816$ and $0.6817\,\mathrm{kpc}$, so the quarter-percent flat regime begins between $2.04$ and $2.05\,\mathrm{kpc}$. Its asymptotic speed lies between $50.2$ and $50.3\,\mathrm{km/s}$. At the surface of the ball the speed lies between $55.2$ and $55.3\,\mathrm{km/s}$; by $2\,\mathrm{kpc}$ it has fallen to between $50.38$ and $50.39\,\mathrm{km/s}$, and at $4\,\mathrm{kpc}$ it is back between $50.2$ and $50.3\,\mathrm{km/s}$. From that surface to $4\,\mathrm{kpc}$ the radius grows by a factor of five, and the specific angular momentum of the circular orbit by a factor between $4.53$ and $4.56$. On the radial lattice the antipode would lie between $2.141$ and $2.142\,\mathrm{kpc}$, inside that flat stretch.

\section{What the sky can decide}
\label{sec:sky}

Neighbouring patches may overlap, each faithful near its own centre. If galaxies and clusters are such patches, the excess of the observed acceleration over the baryonic Newtonian value is the dual image lattice, not an unseen mass, and a direct search has nothing particle-like to find. Three features separate this reading from dark halos and from fitted interpolations alike. Outside the baryons the curve approaches $v_\infty=(GM_{\text{b}}a_0)^{1/4}$ from above and is flat to a quarter of a percent past $3R$; it neither declines like a Keplerian tail nor rises again. The baryonic Tully--Fisher relation $v_\infty^4=GM_{\text{b}}a_0$ is exact, with no dependence on surface brightness or size. And the relation \eqref{eq:dual-rar} has a definite transition that lies lower than the empirical one. At $g_{\text{N}}=a_0$ the enhancement is $1.31$--$1.32$, where the fit of McGaugh, Lelli and Schombert gives about $1.58$. In the galaxies above the dual curve runs up to a tenth below that fit, and the two converge only slowly toward the common $v_\infty$. In star-dominated disks a mass-to-light ratio higher by roughly a fifth absorbs the difference. In gas-dominated dwarfs, whose baryonic curve is fixed by 21-cm data, nothing does, and they decide between the lattice and the fitted relation. The alternation of attraction and repulsion in strong fields, ordered by the interference of equal rapidities, is another regime. It is not the sign of $g$ on this chart. The sense of the disks, read against the one screw, asks which of the two families this universe carries.

\begin{thebibliography}{99}
\raggedright

\bibitem{dst-main}
Anonymous,
``Double Spacetime Theory: Gravity as Torsion between Usual and Dual Spacetime'' (2025).
\url{https://github.com/hypernumbernet/double-spacetime-doc/blob/main/double-spacetime-theory.pdf}.

\end{thebibliography}

\end{document}
